\section*{Capítulo \ref{ch:b3:remote-sensing-gis} --- Sensoriamento remoto e SIG}

\begin{solution}{exo:b3:remote-sensing-gis:1}
Campo: $(0.42 - 0.06)/(0.42 + 0.06) = 0.75$, vegetação verde densa, uma lavoura
em plena folhagem. Campo arado: $(0.28 - 0.21)/(0.28 + 0.21) = 0.14$, quase nu,
um solo que clareia levemente em direção ao infravermelho.
\end{solution}

\begin{solution}{exo:b3:remote-sensing-gis:2}
(a)~$185000/30 \approx 6170$ pixels; $185000/15 \approx 12300$ na
banda pancromática. (b)~$2^{12} = 4096$ níveis; $2^8 = 256$. (c)~16 dias
sozinho; com o Landsat~9 na mesma órbita, defasado de 8 dias, 8 dias.
\end{solution}

\begin{solution}{exo:b3:remote-sensing-gis:3}
(a)~Uma franja vale $\lambda/2 = \qty{27.7}{mm}$; seis franjas
correspondem a \qty{166}{mm} de variação na \omterm{def:b3:remote-sensing-gis:insar}{linha de visada}, aqui afastando-se do satélite, já que a bacia subside. 
(b)~$56.6/2 = \qty{28.3}{mm}$.
\end{solution}

\begin{solution}{exo:b3:remote-sensing-gis:4}
(a)~$t = 2d/c = 2\times\num{5e5}/\num{3.00e8} = \qty{3.34}{ms}$.
(b)~$c\,\delta t/2 = \num{3.00e8}\times\num{8e-10}/2 = \qty{0.12}{m}$; uma
superfície amostrada por muitos fótons é localizada com muito mais precisão, a
precisão caindo com o inverso da raiz quadrada do número deles.
\end{solution}

\begin{solution}{exo:b3:remote-sensing-gis:5}
(a)~No SI, $L = \qty{9.0e6}{W.m^{-2}.sr^{-1}.m^{-1}}$. Com $c_1 =
2hc^2 = \qty{1.191e-16}{W.m^2.sr^{-1}}$, $c_2 = \qty{0.014388}{m.K}$ e
$\lambda^5 = (\num{1.09e-5})^5 = \num{1.539e-25}$, $c_1/(\lambda^5L) = 86.0$,
e $T_b = c_2/(\lambda\ln 87.0) = \qty{295.6}{K}$. (b)~A temperatura cinética
resolve $0.98\,B_\lambda(T) = L$, isto é, $B_\lambda(T) = L/0.98$: a mesma
fórmula dá $T = \qty{296.9}{K}$, \qty{1.3}{K} a mais.
\end{solution}

\begin{solution}{exo:b3:remote-sensing-gis:6}
(a)~$R = 693/\cos\ang{35} = \qty{846}{km}$. (b)~Antena real: $R\lambda/L =
\num{8.46e5}\times0.0555/12.3 = \qty{3.8}{km}$; abertura sintética: $L/2 =
\qty{6.2}{m}$, cerca de 600 vezes mais fina. (c)~A abertura sintética é o
comprimento da trajetória ao longo do qual um ponto permanece no feixe, $R\lambda/L =
\qty{3.8}{km}$.
\end{solution}

\begin{solution}{exo:b3:remote-sensing-gis:7}
$\partial z/\partial x = (418 - 405)/60 = 0.217$ (leste menos oeste),
$\partial z/\partial y = (418 - 401)/60 = 0.283$ (norte menos sul);
$|\nabla z| = 0.357$, uma declividade de $\arctan 0.357 = \ang{19.6}$. O terreno
sobe para nordeste, logo sua face está voltada para sudoeste: $A =
\operatorname{atan2}(-0.217, -0.283) = \ang{217}$. Sombreamento: $\cos\ang{45}
\cos\ang{19.6} + \sin\ang{45}\sin\ang{19.6}\cos(\ang{315} - \ang{217}) =
0.666 - 0.033 = 0.63$, um pouco mais escuro que o terreno plano (0.71): a
encosta se afasta levemente de uma luz no noroeste.
\end{solution}

\begin{solution}{exo:b3:remote-sensing-gis:8}
Totais das linhas 52, 48, 38, 22; totais das colunas 54, 46, 38, 22; diagonal 133 de
160: exatidão global \qty{83.1}{\percent}. Exatidões do usuário: água
$48/52 = 0.92$, floresta $37/48 = 0.77$, cultivos $29/38 = 0.76$, cidade
$19/22 = 0.86$. Do produtor: água $48/54 = 0.89$, floresta $37/46 = 0.80$,
cultivos $29/38 = 0.76$, cidade $19/22 = 0.86$. $p_e = (52\cdot54 + 48\cdot46 +
38\cdot38 + 22\cdot22)/160^2 = 0.271$, $\kappa = (0.831 - 0.271)/(1 -
0.271) = 0.77$. Os cultivos são os mais frequentemente não identificados: 9 dos 38 pixels verdadeiros de cultivo
foram para outro lugar, principalmente para floresta.
\end{solution}

\begin{solution}{exo:b3:remote-sensing-gis:9}
(a)~$L = 0.04\,B_\lambda(1443) + 0.96\,B_\lambda(290)$ em
\qty{10.9}{\micro\metre}, \qty{28.6}{W.m^{-2}.sr^{-1}.\micro\metre^{-1}},
cuja \omterm{def:b3:remote-sensing-gis:brightness}{temperatura de brilho} é \qty{396}{K}, contra uma temperatura média de
$0.04\times1443 + 0.96\times290 = \qty{336}{K}$. (b)~Em
\qty{2.2}{\micro\metre}, $B_\lambda(1443)/B_\lambda(290) = \num{6.8e7}$:
a lava domina completamente o pixel, e sua \omterm{def:b3:remote-sensing-gis:brightness}{temperatura de brilho} é
\qty{845}{K}. A banda do infravermelho de ondas curtas detecta pontos quentes pequenos
muito melhor; a banda térmica, sensível ao pixel inteiro, é a que
mede o terreno ao redor deles.
\end{solution}

\begin{solution}{exo:b3:remote-sensing-gis:10}
Com o satélite a leste, $d_1 = u_E\sin\theta + u_z\cos\theta$; a
oeste, $d_2 = -u_E\sin\theta + u_z\cos\theta$. Somando e subtraindo,
$u_z = (d_1 + d_2)/(2\cos\theta) = 102/1.638 = \qty{62.3}{mm}$ e $u_E =
(d_1 - d_2)/(2\sin\theta) = 26/1.147 = \qty{22.7}{mm}$. Uma passagem dá uma
equação para duas incógnitas: qualquer mistura de soerguimento e movimento
para leste com a mesma projeção na \omterm{def:b3:remote-sensing-gis:insar}{linha de visada} produz as mesmas franjas.
\end{solution}

\begin{solution}{exo:b3:remote-sensing-gis:11}
(a)~Pelo \cref{thm:b3:remote-sensing-gis:fringe}, $\Delta\phi =
4\pi\Delta r/\lambda < \pi$ exige $\Delta r < \lambda/4 =
\qty{13.9}{mm}$. (b)~$u_z = k z_0(r^2 + z_0^2)^{-3/2}$ com $k = (1 -
\nu)\Delta V/\pi$; $\mathrm du_z/\mathrm dr = -3kz_0r(r^2 + z_0^2)^{-5/2}$,
cuja derivada se anula onde $(r^2 + z_0^2) = 5r^2$, em $r = z_0/2$. Aí
\[
\left|\frac{\mathrm du_z}{\mathrm dr}\right| = \frac{3k}{2z_0^3}\Big(\frac54\Big)^{-5/2}
= \frac{0.859\,k}{z_0^3} = \frac{0.859\times\num{2.39e6}}{\num{6.4e10}} = \num{3.2e-5},
\]
\qty{32}{mm} por quilômetro. (c)~$13.9/32 = \qty{0.43}{km}$: pixels menores
que cerca de \qty{430}{m}, facilmente atendidos pelo Sentinel-1.
\end{solution}

\begin{solution}{exo:b3:remote-sensing-gis:12}
(a)~Para $N, R \ge 0$ não ambos nulos, $|N - R| \le N + R$, logo $|\mathrm{NDVI}|
\le 1$; substituir $N$ e $R$ por $aN$ e $aR$ multiplica o numerador e
o denominador por $a$. (b)~Sem correção: $(0.502 - 0.051)/(0.502 + 0.051) =
0.816$; com correção: $(0.512 - 0.081)/(0.512 + 0.081) = 0.727$. (c)~A
radiância do caminho reduz o índice da vegetação verde, cerca de um décimo aqui, e
em uma quantidade que muda com a neblina: um mapa de imagens não corrigidas
mostraria o tempo tanto quanto a vegetação, e as comparações entre datas
não teriam sentido.
\end{solution}

\begin{solution}{pb:b3:remote-sensing-gis:1}
\textbf{1.} Floresta $(0.38 - 0.04)/0.42 = 0.81$; fluxo de 1890 $0.23/0.37 =
0.62$; de 1960 $0.10/0.32 = 0.31$; de 2005 $0.02/0.26 = 0.08$.

\textbf{2.} 1890, depois 1960, depois 2005: o fluxo mais antigo é o mais verde,
na ordem das idades; ele ainda não alcançou a floresta.

\textbf{3.} A rocha nua reflete um pouco mais no infravermelho próximo do que no
vermelho, como a areia marrom da figura do capítulo; líquens e as primeiras
plantas acrescentam um pouco: 0.08 é o valor de um terreno quase nu.

\textbf{4.} Uma encosta voltada para longe da luz escurece as duas bandas na mesma
proporção; a razão cancela esse efeito, enquanto a refletância no infravermelho próximo sozinha
confundiria uma floresta sombreada com vegetação esparsa.

\textbf{5.} Pela \cref{prop:b3:remote-sensing-gis:planck-inversion} com
$\lambda = \qty{10.9}{\micro\metre}$: o platô \qty{291.1}{K}, o
pixel da fumarola \qty{302.3}{K}.

\textbf{6.} $0.97\,B_\lambda(T) = \num{8.40e6}$ resulta em $T = \qty{293.0}{K}$.

\textbf{7.} $B_\lambda(\qty{370}{K}) =
\qty{22.48}{W.m^{-2}.sr^{-1}.\micro\metre^{-1}}$; $f = (9.95 - 8.40)/(22.48
- 8.40) = 0.110$, cerca de \qty{1100}{m^2} do pixel de \qty{e4}{m^2}.

\textbf{8.} A radiância é uma função convexa da temperatura
(\cref{prop:b3:remote-sensing-gis:mixed}): a parte quente acrescenta mais radiância
do que sua parcela na temperatura, de modo que o pixel parece mais quente do que sua média,
$0.110\times370 + 0.890\times291.1 = \qty{299.8}{K}$ frente a
\qty{302.3}{K}.

\textbf{9.} A fase de cada eco depende do caminho exato; uma grande
mudança de ponto de vista altera a soma dos ecos dos espalhadores, e o
\omterm{def:b3:remote-sensing-gis:insar}{interferograma} perde sua coerência e ganha uma fase devida ao relevo.
$48/12 = 4$ ciclos de um satélite.

\textbf{10.} $5.4\times27.73 = \qty{149.7}{mm}$ em direção ao satélite.

\textbf{11.} $u_z = 149.7/\cos\ang{35} = \qty{182.8}{mm}$.

\textbf{12.} O solo se move para fora, além de para cima; no lado leste o
movimento para fora é em direção ao satélite e soma-se ao soerguimento; no lado oeste
ele subtrai (\cref{prop:b3:remote-sensing-gis:mogi-los}). A variação
cai mais depressa para oeste, onde as franjas estão concentradas; o padrão está
alongado para leste.

\textbf{13.} Nessa linha, o movimento horizontal é norte--sul,
perpendicular à \omterm{def:b3:remote-sensing-gis:insar}{linha de visada}, de modo que o radar vê somente $u_z\cos\theta$, com o formato do
soerguimento. O soerguimento cai à metade em $r = 0.766\,z_0$, logo $z_0 =
3.4/0.766 = \qty{4.44}{km}$.

\textbf{14.} Metade de 5.4: 2.7 franjas.

\textbf{15.} Uma deformação é a mesma nos dois \omterm{def:b3:remote-sensing-gis:insar}{interferogramas} após a normalização pelos intervalos de tempo,
ou continua crescendo de par para par; um padrão atmosférico
muda de data para data e costuma acompanhar o
relevo. Um segundo par que mostra o mesmo padrão de olho de boi confirma a
deformação.

\textbf{16.} $\Delta V = \pi u_z(0)z_0^2/(1 - \nu) = \pi\times0.1828\times
4436^2/0.75 = \qty{1.51e7}{m^3}$.

\textbf{17.} $\num{1.51e7}/48 = \qty{3.1e5}{m^3}$ por dia, \qty{1.1e8}{m^3}
por ano, se a taxa se mantivesse.

\textbf{18.} $u_r \propto r/(r^2 + z_0^2)^{3/2}$ é maior onde sua
derivada se anula, $r^2 + z_0^2 = 3r^2$, $r = z_0/\sqrt2 = \qty{3.14}{km}$;
ali $u_r = \qty{70}{mm}$, \qty{38}{\percent} do soerguimento central.

\textbf{19.} Raio \qty{1.48}{km}, volume $\frac43\pi(1.48)^3 =
\qty{13.6}{km^3}$; a variação é $\num{1.51e7}/\num{1.36e10} =
\qty{0.11}{\percent}$ do volume.

\textbf{20.} Não: um décimo de franja é \qty{2}{\percent} da contagem, e
a profundidade, lida a partir de uma meia-largura de \qty{3.4}{km}, talvez até
\qty{0.1}{km}, entra elevada ao quadrado, \qty{6}{\percent}; a correção de \qty{3}{\percent}
está dentro dessa incerteza.

\textbf{21.} O aquecimento do campo de fumarolas (Parte~II) ao longo do tempo, um
aumento de pequenos terremotos sob o abaulamento, uma mudança no gás
liberado, e o soerguimento medido por posicionamento por satélite ou nivelamento em
alguns pontos do flanco.

\textbf{22.} O reservatório sob Mount Halvren cresceu em
\textbf{15~milhões~de~\unit[text-series-to-math=true]{m^3}} ($\num{1.51e7}\,\unit{m^3}$).
\end{solution}
